\documentclass[11pt,a4paper]{article}
\usepackage[margin=1in]{geometry}
\usepackage{amsmath,amssymb}
\usepackage{adjustbox}
\usepackage{pdflscape}
\usepackage[colorlinks=true,allcolors=blue]{hyperref}
\setlength{\parindent}{0pt}
\setlength{\parskip}{0.6em}
\emergencystretch=3em

\title{The overlap Dirac determinant is not link-reflection positive: a certified witness on an $8\times 1$ lattice%
\thanks{Multi-DAC. Authors in alphabetical order. Clawd and Wren are AIs. How this work was done and checked is described in the README of
the repository \texttt{Multi-DAC/releases}, where the scripts for every computed figure are filed with this
paper.}}
\author{Clawd \and Clayton Iggulden-Schnell \and Wren}
\date{25 November 2026}

\begin{document}
\maketitle

\begin{abstract}
Reflection positivity is the property of a Euclidean lattice measure that yields a Hilbert space and a positive transfer matrix. For the overlap Dirac operator it is proved for free fermions (Kikukawa and Usui) and open in a gauge background. We show that it fails, for the reflection through a plane of links, for the measure $\det D_{\rm ov}[U]\,e^{-\beta S_G[U]}\,dU$ of two-dimensional $U(1)$ lattice gauge theory on an $8\times 1$ lattice, with Wilson kernel ($r=1$, kernel mass $m=0.7$) and the Wilson plaquette action. For every $\beta\ge 0$ there is a smooth, real, gauge-invariant observable $F$ with $\langle(\Theta F)F\rangle_\beta<0$, and one $F$ serves every $\beta$ in any bounded interval. The witness is the $2\times 2$ Gram matrix $K$ of two gauge fields that differ in one spatial link: $v^TKv\in[-8.2647\times10^{-9},\,-8.2322\times10^{-9}]$ for $v=(1,-1)$, enclosed by two interval-arithmetic computations that share no arithmetic code. The fields are smooth (every plaquette has $|1-U_p|\le 0.1997$), so the failure is not a rough-field effect and does not come from an admissibility cut. Applied to the Wilson determinant, which is reflection positive by theorem, the same construction gives a positive semidefinite Gram matrix.
\end{abstract}


\section*{1. Introduction}
A lattice measure is reflection positive if $\langle(\Theta F)F\rangle\ge 0$ for every observable $F$ supported on one side of a reflection plane, where $\Theta$ reflects and complex-conjugates. Reflection positivity is what turns a Euclidean lattice theory into a quantum theory: it gives a Hilbert space and a positive transfer matrix. For Wilson fermions with $r=1$ and the Wilson plaquette action it was proved by Osterwalder and Seiler [1], for reflection through a plane of links; Jaffe and Janssens [2] give a general criterion and recover the Wilson case for every real mass.

The overlap operator of Neuberger [3] has an exact lattice chiral symmetry. Whether it is reflection positive in a gauge background has been open. Kikukawa and Usui [4] proved reflection positivity for free overlap fermions at equal lattice spacings, and left the gauge case for future work. Mandula [5] wrote that overlap fermions do not satisfy reflection positivity, inferring this from branch points of the free kernel in the limit of vanishing time spacing; no expectation value was shown to be negative. Creutz [6] showed that a single-plaquette gauge weight that vanishes on an arc, as an admissibility cut on the plaquettes requires, cannot give a positive transfer matrix; such cuts are used to guarantee that the overlap operator is well defined. Bietenholz et al. [7] accept this for their topology-conserving gauge actions and check a weaker property, site reflection positivity. Kikukawa [8] states that the overlap construction of a chiral gauge theory is reflection positive in the weak gauge-coupling limit, and that at finite coupling reflection positivity is lost already in the gauge action, because of the admissibility condition. Lu and Shao [9] show that the free overlap Weyl fermion admits no reflection of the required kind, a different statement about a different object.

Here the gauge action is the plain Wilson plaquette action, with no cut, and the failure is in the determinant. We give two smooth gauge fields on an $8\times 1$ lattice whose $2\times2$ reflection Gram matrix has a negative eigenvalue, and turn it into a smooth gauge-invariant observable with $\langle(\Theta F)F\rangle_\beta<0$ at every $\beta$. The negative value is certified by interval arithmetic. We know of no earlier example of a negative reflection expectation value for the overlap determinant in a gauge background.


\section*{2. Setting and result}
{\bfseries\boldmath Lattice and fields.} Time sites are $t=-3,\dots,4$ ($2L=8$, $L=4$), with one spatial site, periodic in space. The gauge group is $U(1)$. The temporal link $U_0(t)$ runs from $t$ to $t+1$, and $U_0(4)$ runs from $4$ to $-3$. The spatial link $U_1(t)$ runs from the site at time $t$ around the spatial circle to itself. Fermions are antiperiodic in time.

{\bfseries\boldmath Operators.} With $\gamma_0=\sigma_1$, $\gamma_1=\sigma_2$ and $\gamma_5=\sigma_3$, the Wilson operator with $r=1$ is


\begin{equation*}
(D_W\psi)(x)=2\psi(x)-\frac12\sum_{\mu=0,1}\Big[(1-\gamma_\mu)U_\mu(x)\psi(x+\hat\mu)+(1+\gamma_\mu)\overline{U_\mu(x-\hat\mu)}\psi(x-\hat\mu)\Big],
\end{equation*}


with a factor $-1$ on the temporal hop across the link $4\to-3$. The kernel is $X=D_W-m$ with $m=0.7$, $H=X^\dagger X$, and the overlap operator is


\begin{equation*}
D_{\rm ov}=\tfrac12\big(1+XH^{-1/2}\big),
\end{equation*}


defined wherever $\det X\neq0$. It acts on $\mathbb{C}^{16}$ (8 sites, 2 spinor components).

{\bfseries\boldmath Measure.} The gauge action is $S_G[U]=\sum_p(1-\mathrm{Re}\,U_p)$, and


\begin{equation*}
d\mu_\beta=Z^{-1}\det D_{\rm ov}[U]\,e^{-\beta S_G[U]}\,dU,\qquad \beta\ge0,
\end{equation*}


with $dU$ the Haar measure on the 16 links. Lemma 5 shows $0\le\det D_{\rm ov}\le1$ almost everywhere and $0<Z<\infty$.

{\bfseries\boldmath Reflection.} $\theta(t)=1-t$ fixes the plane between $t=0$ and $t=1$. Spatial links are copied, $(\theta U)_1(1-t)=U_1(t)$, and temporal links are reversed and conjugated, $(\theta U)_0(-t)=\overline{U_0(t)}$. The positive half is $\Lambda_+=\{1,2,3,4\}$, and the links $U_0(0)$ and $U_0(4)$ cross the plane. For a function $F$ of the links inside $\Lambda_+$, $(\Theta F)(U)=\overline{F(\theta U)}$. The measure is link-reflection positive if $\langle(\Theta F)F\rangle_\beta\ge0$ for all such $F$.

{\bfseries\boldmath Theorem 1.} For every $\beta\ge0$ there is a smooth, real, gauge-invariant function $F_\beta$ of the spatial links $U_1(1),\dots,U_1(4)$ with $\langle(\Theta F_\beta)F_\beta\rangle_\beta<0$. For every $B>0$ one such $F$ serves every $\beta\in[0,B]$. So $\mu_\beta$ is not link-reflection positive at any $\beta\ge0$.

{\bfseries\boldmath The witness.} Let $\epsilon=0.2$. The field $A_1$ has $U_1(2)=e^{i\epsilon}$ and $A_2$ has $U_1(3)=e^{i\epsilon}$; every other link of $\Lambda_+$ is $1$ in both. Write $[\theta A_i,\Phi,A_j]$ for the configuration whose spatial links are those of $\theta A_i$ on $t\le0$ and of $A_j$ on $t\ge1$, and whose temporal links have total holonomy $\prod_tU_0(t)=e^{i\Phi}$ (by Lemma 3(b) only the holonomy matters). Put


\begin{equation*}
K_{ij}=\frac1{2\pi}\int_0^{2\pi}d\Phi\,\det D_{\rm ov}[\theta A_i,\Phi,A_j].
\end{equation*}


{\bfseries\boldmath Proposition 2 (certified).} With $v=(1,-1)$,


\begin{equation*}
v^TKv=K_{11}+K_{22}-K_{12}-K_{21}\in[-8.264675148\times10^{-9},\,-8.232272695\times10^{-9}].
\end{equation*}


The centres are $K_{11}=K_{22}=3.784508295226\times10^{-4}$ and $K_{12}=K_{21}=3.784549537596\times10^{-4}$.

{\bfseries\boldmath What is not claimed.} One $F$ for all $\beta\ge0$ is not claimed: the smearing radius in the proof shrinks as $\beta$ grows. The theorem is for one lattice, one kernel mass and one impurity angle. Larger spatial extent, other $L$ and other $m$ are not certified here, nor is anything about four dimensions, non-abelian groups or the continuum limit.

{\bfseries\boldmath Fermion boundary condition.} The antiperiodic sign on the wrap link is the shift $\Phi\to\Phi+\pi$ of the holonomy, which the $\Phi$ integral absorbs. So $K$, and Theorem 1, are the same for periodic fermions.


\section*{3. Proof}

\subsection*{3.1 Reduction to one holonomy}
{\bfseries\boldmath Lemma 3.}
(a) Each spatial link $U_1(t)$ is gauge invariant, so every function of $U_1(1),\dots,U_1(4)$ is gauge invariant and supported in $\Lambda_+$.
(b) Every configuration of the temporal links is gauge-equivalent, without changing any $U_1$, to $U_0(t)=e^{i\Phi/8}$ on every temporal link, where $e^{i\Phi}=\prod_tU_0(t)$; Haar measure on the temporal links pushes forward to $d\Phi/2\pi$.
(c) Every plaquette is $U_p(t)=U_1(t+1)\overline{U_1(t)}$ (with $U_1(5)\equiv U_1(-3)$), so $S_G$ does not depend on the temporal links.
(d) Write $u_a(t)=(\theta U)_1(t)=U_1(1-t)$ and $u_b(t)=U_1(t)$ for $t=1,\dots,4$, and $S_+(u)=\sum_{t=1}^{3}(1-\mathrm{Re}\,u(t+1)\overline{u(t)})$. Then


\begin{equation*}
S_G=S_+(u_a)+S_+(u_b)+\big(1-\mathrm{Re}\,u_b(1)\overline{u_a(1)}\big)+\big(1-\mathrm{Re}\,u_a(4)\overline{u_b(4)}\big).
\end{equation*}


\emph{Proof.} (a) With one spatial site, a gauge transformation $g$ acts by $U_1(t)\mapsto g(t)U_1(t)\overline{g(t)}=U_1(t)$. (b) The gauge transformation $g(t)=\prod_{s<t}U_0(s)\,e^{-i\Phi(t+3)/8}$ (product from $s=-3$) moves every temporal link to $e^{i\Phi/8}$ and leaves the abelian spatial links unchanged. The holonomy of independent Haar-distributed links is Haar distributed. (c) The two temporal links of each plaquette cancel. (d) The plaquettes between $t$ and $t+1$ for $t=1,2,3$ give $S_+(u_b)$, those for $t=-3,-2,-1$ give $S_+(u_a)$, and the rows $0\to1$ and $4\to-3$ are the two crossing terms. $\square$

For a real $F$ of $U_1(1),\dots,U_1(4)$ we have $(\Theta F)(U)=F(u_a)$. Put


\begin{equation*}
\mathcal K_\beta(u_a,u_b)=\frac1{2\pi}\int_0^{2\pi}d\Phi\,\det D_{\rm ov}[\theta u_a,\Phi,u_b]\,e^{-\beta S_G[\theta u_a,u_b]}.
\end{equation*}


Then by Lemma 3


\begin{equation*}
Z\langle(\Theta F)F\rangle_\beta=\int\!\!\int F(u_a)F(u_b)\,\mathcal K_\beta(u_a,u_b)\,du_a\,du_b. \tag{1}
\end{equation*}


{\bfseries\boldmath At the witness.} Both $A_1$ and $A_2$ have $U_1(1)=U_1(4)=1$, so both crossing terms in Lemma 3(d) vanish for every pair $(A_i,A_j)$, and $S_+(A_1)=S_+(A_2)=2(1-\cos\epsilon)$. Hence


\begin{equation*}
\mathcal K_\beta(A_i,A_j)=c(\beta)K_{ij},\qquad c(\beta)=e^{-4\beta(1-\cos\epsilon)}>0. \tag{2}
\end{equation*}


This uses the shared boundary links. For a family in which $U_1(1)$ or $U_1(4)$ varies, the crossing terms multiply $K$ entrywise by a matrix that does not factorise, and (2) fails.

At one spatial site the gauge weight by itself is reflection positive: its crossing factor is $e^{\beta\cos(\alpha-\alpha')}$ for each crossing row, a positive-definite kernel on $U(1)$ because its Fourier coefficients are the Bessel values $I_n(\beta)\ge0$. Any failure of reflection positivity therefore comes from the determinant.


\subsection*{3.2 The certificate for Proposition 2}
Fix one of the four pairs $(i,j)$, and let $X(\Phi)$ be the kernel at $[\theta A_i,\Phi,A_j]$ with every temporal link equal to $e^{i\Phi/8}$.

{\bfseries\boldmath Bounds on the kernel.} The temporal hops are $-(S\otimes P_-)e^{i\Phi/8}-(S^\dagger\otimes P_+)e^{-i\Phi/8}$, with $S$ the signed unitary shift and $P_\mp=(1\mp\gamma_0)/2$. Since $P_-P_+=0$, the norm of $c_1(S\otimes P_-)+c_2(S^\dagger\otimes P_+)$ is $\max(|c_1|,|c_2|)$. The spatial hop at time $t$ is $-(P^{(1)}_-U_1(t)+P^{(1)}_+\overline{U_1(t)})$, with complementary projectors $P^{(1)}_\mp=(1\mp\gamma_1)/2$, and it is block-diagonal in $t$. Hence the temporal and spatial hops each have norm $1$, $\|X(\Phi)\|\le M_X=|2-m|+1+1=3.3$, and $\|dX/d\Phi\|=1/8$. By Weyl's inequality $\sigma_{\min}(X(\Phi))$ is $\tfrac18$-Lipschitz in $\Phi$.

{\bfseries\boldmath Part 1: a gap for every $\Phi$.} At the $512$ angles $\Phi_g=2\pi g/512$, an interval Cholesky factorisation of the real $32\times32$ form of $H-0.6848^2$ certifies $\sigma_{\min}(X(\Phi_g))>0.6848$. The Lipschitz bound then gives $\sigma_{\min}(X(\Phi))\ge0.6840330$ for every $\Phi$, so $\lambda_{\min}(H(\Phi))\ge\lambda_0=0.4679011$, for all four pairs.

{\bfseries\boldmath Part 2: a strip, and the quadrature error.} Continue to $z=\Phi+iy$ with $|y|\le\eta=0.47$, and put $H(z)=X(\bar z)^\dagger X(z)$, which is holomorphic and equals $H(\Phi)$ on the real axis. Then $\|X(z)-X(\Phi)\|\le\delta=e^{\eta/8}-1\le0.0606$ and $\|H(z)-H(\Phi)\|\le\Delta=2M_X\delta+\delta^2\le0.4031$. So the numerical range of $H(z)$ has real part at least $c_0=\lambda_0-\Delta\ge0.06487$. For such an $H$, $\|(t+H)^{-1}\|\le(t+c_0)^{-1}$ for $t\ge0$, and


\begin{equation*}
H(z)^{-1/2}=\frac1\pi\int_0^\infty t^{-1/2}(t+H(z))^{-1}\,dt
\end{equation*}


is holomorphic in $z$, extends $H(\Phi)^{-1/2}$, and has norm at most $c_0^{-1/2}$. Hence $D(z)=\tfrac12(1+X(z)H(z)^{-1/2})$ is holomorphic on the strip, with $\|D(z)\|\le M_D=\tfrac12(1+(M_X+\delta)c_0^{-1/2})\le7.0970$ and $|\det D(z)|\le M_D^{16}\le4.1414\times10^{13}$.

The function $\det D(z)$ has period $2\pi$: the shift $\Phi\to\Phi+2\pi$ multiplies every temporal link by $e^{2\pi i/8}$, which the $z$-independent gauge transformation $\mathrm{diag}(e^{2\pi it/8})$ undoes. For a $2\pi$-periodic $g$, holomorphic on a neighbourhood of $|\mathrm{Im}\,z|\le\eta$ with $|g|\le M$ there, the Fourier coefficients obey $|\hat g_k|\le Me^{-|k|\eta}$ (shift the contour), and the $n$-point trapezoid rule has error $|\sum_{j\ne0}\hat g_{jn}|\le 2M/(e^{n\eta}-1)$. With $n=124$ this is at most $Q=4.0504\times10^{-12}$ per entry, and $4Q\le1.621\times10^{-11}$ on $v^TKv$.

{\bfseries\boldmath Part 3: the determinant at each node.} At each node $\Phi_k=2\pi k/124$ and for each pair: let $M_0$ be an exactly Hermitian high-precision approximation of $H^{1/2}$; bound $r\ge\|M_0^2-H\|_F$ in interval arithmetic; and certify $M_0\succeq aI$ by an interval Cholesky factorisation of $M_0-aI$, with $a^2>r$. Then $H\succeq(a^2-r)I$, and the integral representation gives


\begin{equation*}
\|H^{-1/2}-M_0^{-1}\|\le\frac1\pi\int_0^\infty\frac{t^{-1/2}\,r\,dt}{(t+a^2-r)(t+a^2)}=\rho=\frac1{\sqrt{a^2-r}}-\frac1a.
\end{equation*}


Write $D_{\rm ov}=D_0+E$ with $D_0=\tfrac12(1+XM_0^{-1})$, so that $\det D_0=\det(M_0+X)/(2^{16}\det M_0)$ is computed in interval arithmetic, $\|D_0\|\le\tfrac12(1+M_X/a)$ and $\|E\|\le M_X\rho/2$. Since $\tfrac{d}{ds}\det(D_0+sE)=\mathrm{tr}(\mathrm{adj}(D_0+sE)E)$ and the trace norm of $\mathrm{adj}\,C$ is at most $N\|C\|^{N-1}$,


\begin{equation*}
|\det(D_0+E)-\det D_0|\le N\|E\|\,(\|D_0\|+\|E\|)^{N-1},\qquad N=16.
\end{equation*}


Across all nodes and pairs this error is at most $2.7\times10^{-28}$.

{\bfseries\boldmath Result.} $v^TKv$ lies in the trapezoid sum of the enclosed $\det D_0$ values, widened by the node errors and by $[-4Q,4Q]$. Its real part lies in $[-8.264675148\times10^{-9},-8.232272695\times10^{-9}]$; the total radius is $1.62\times10^{-11}$, almost all of it quadrature. The imaginary part lies in $[-1.62\times10^{-11},1.62\times10^{-11}]$, consistent with $K$ being real.

{\bfseries\boldmath Second computation.} The interval exponential, sine and cosine of the library used above (mpmath 1.3.0) are not formally outward-rounded. A second program, written separately, avoids them: it computes $\pi$ by Machin's series with a tail bound, and $\exp$, $\cos$ and $\sin$ by Taylor polynomials with Lagrange remainders, using only interval $+,-,\times,\div$. It takes exact rational inputs, uses its own interval Gaussian elimination, and works at 160 bits. It gives $v^TKv\in[-8.26467514775\times10^{-9},-8.23227269469\times10^{-9}]$. The two programs share the definition of the operator and the design of the bounds, and no code. Section 4 tests the definition of the operator.


\subsection*{3.3 From the Gram matrix to an observable}
{\bfseries\boldmath Lemma 4 (continuity).} Let $0<\delta<0.684$, and let $N_\delta\subset U(1)^4$ be the union of the closed balls $\max_t|u(t)-A_i(t)|\le\delta$ around $A_1$ and $A_2$. For $u_a,u_b\in N_\delta$ and every $\Phi$, $H[\theta u_a,\Phi,u_b]\succeq(0.6840-\delta)^2>0$. So $\det D_{\rm ov}$ is real-analytic there, and $\mathcal K_\beta(u_a,u_b)$ is jointly continuous on $[0,B]\times N_\delta\times N_\delta$ for every $B>0$.

\emph{Proof.} By the form of the spatial hop in Section 3.2, moving the spatial links by at most $\delta$ moves $X$ by at most $\delta$ in norm, and Weyl's inequality with Part 1 gives the bound. $\square$

{\bfseries\boldmath Smearing.} Let $\varphi_r\ge0$ be a smooth function on $U(1)^4$ with $\int\varphi_r=1$, supported in the ball of radius $r<\delta$ at the identity, and put


\begin{equation*}
F_r(u)=\varphi_r(uA_1^{-1})-\varphi_r(uA_2^{-1}).
\end{equation*}


$F_r$ is real, smooth, and a function of $U_1(1),\dots,U_1(4)$ alone, hence gauge invariant by Lemma 3(a). By (1),


\begin{equation*}
Z\langle(\Theta F_r)F_r\rangle_\beta=\sum_{i,j}v_iv_j\int\!\!\int\varphi_r(u_aA_i^{-1})\varphi_r(u_bA_j^{-1})\,\mathcal K_\beta(u_a,u_b)\,du_a\,du_b.
\end{equation*}


$\mathcal K_\beta$ is uniformly continuous on the compact set $[0,B]\times N_\delta\times N_\delta$ (Lemma 4), with modulus of continuity $\omega$. Since $\varphi_r\ge0$ has mass $1$, the right side differs from $\sum_{ij}v_iv_j\mathcal K_\beta(A_i,A_j)=c(\beta)\,v^TKv$ by at most $4\omega(r)$, uniformly in $\beta\in[0,B]$. By (2) and Proposition 2, $c(\beta)v^TKv\le c(B)v^TKv<0$ on $[0,B]$. So one $r>0$ makes $Z\langle(\Theta F_r)F_r\rangle_\beta<0$ for every $\beta\in[0,B]$.

{\bfseries\boldmath Lemma 5.} $0\le\det D_{\rm ov}\le1$ almost everywhere, and $0<Z<\infty$ for every $\beta\ge0$.

\emph{Proof.} The set $\{\det X=0\}$ is the zero set of $\det H=|\det X|^2$, a trigonometric polynomial in the link angles. It is not identically zero (Part 1), so the set is Haar-null. Off it, $V=XH^{-1/2}$ is unitary, and $\gamma_5X\gamma_5=X^\dagger$ gives $\gamma_5V\gamma_5=V^\dagger$. So the characteristic polynomial of $V$ has real coefficients: non-real eigenvalues come in conjugate pairs, each pair contributing $|(1+\lambda)/2|^2\in[0,1]$, and the real eigenvalues $\pm1$ contribute $1$ or $0$. Hence $0\le\det D_{\rm ov}=\det\tfrac12(1+V)\le1$, and $Z\le1$. Since $K_{11}>0$, $\det D_{\rm ov}>0$ at some configuration $[\theta A_1,\Phi,A_1]$, and by continuity on an open set, so $Z>0$. $\square$

Dividing by $Z>0$ keeps the sign, which proves Theorem 1. $\blacksquare$


\subsection*{3.4 Remark: the witness is smooth}
Every plaquette of every pair $[\theta A_i,\Phi,A_j]$ satisfies $|1-U_p|\le2\sin(\epsilon/2)=0.19967$: the interior plaquettes carry one link at angle $\epsilon$, and the crossing ones equal $1$. So the witness lies inside every admissible set $|1-U_p|<\epsilon_{\rm adm}$ with $\epsilon_{\rm adm}>0.19967$, and the failure is not caused by rough configurations. Under such a cut the gauge measure itself changes. By Creutz's argument [6], a plaquette weight that is analytic near the identity and vanishes on an arc cannot have all its Fourier coefficients non-negative; for link reflection across a row of plaquettes on $U(1)$, Bochner's theorem makes that the loss of positivity of the crossing kernel. This remark is about where the witness lies, not a theorem about the theory with a cut.


\section*{4. Structure of the witness and a positive control}
{\bfseries\boldmath The Gram matrix is Hankel.} Let $B_s$ ($s=1,\dots,4$) be the field with $U_1(s)=e^{i\epsilon}$ and every other link of $\Lambda_+$ equal to $1$, so $A_1=B_2$ and $A_2=B_3$. In $[\theta B_s,\Phi,B_t]$ the two impurities sit at $1-s$ and $t$ on the ring of 8 time sites, and both carry $+\epsilon$, because $\theta$ copies spatial links without conjugating them. With the holonomy distributed as $e^{i(\Phi+\pi)/8}$ on every temporal link, which absorbs the antiperiodic sign, the free operator is invariant under the $\mathbb{Z}_8$ time translations at every $\Phi$. So $\det D_{\rm ov}[\theta B_s,\Phi,B_t]$ depends only on the separation $d=s+t-1$ modulo 8. Translation by $-d$ maps the impurity set $\{0,d\}$ to $\{0,8-d\}$, so the dependence is through $f(d)=f(8-d)$. Both statements hold pointwise in $\Phi$. After the $\Phi$ integral, $K_{st}=f(s+t-1)$ is a Hankel matrix, and the witness is


\begin{equation*}
v^TKv=2f(3)-2f(4).
\end{equation*}


{\bfseries\boldmath Values (floating point, not certified; 256-point trapezoid rule, $\beta=0$).}

\begin{center}
\begin{adjustbox}{max width=\linewidth}
\begin{tabular}{llll}
\hline
{}{\bfseries\boldmath $d$} & {\bfseries\boldmath $f(d)$, overlap} & {\bfseries\boldmath $f(d)$, $\det(D_W-0.7)$} & {\bfseries\boldmath $f(d)$, $\det(D_W+0.1)$} \\
\hline
{}1 & $3.7168162704528\times10^{-4}$ & $1.0394695346641$ & $6.4161682134288$ \\
{}2 & $3.7807217417675\times10^{-4}$ & $1.0035524990232$ & $6.4029583266392$ \\
{}3 & $3.7845082952260\times10^{-4}$ & $1.0003223032582$ & $6.3952201865087$ \\
{}4 & $3.7845495375956\times10^{-4}$ & $1.0000575572252$ & $6.3926717666078$ \\
\hline
\end{tabular}
\end{adjustbox}
\end{center}

For the overlap, $f$ increases with the ring distance of the two impurities, all the way to the antipode $d=4$; for the Wilson determinant it decreases. The step from $d=3$ to $d=4$ is the negative value, $2f(3)-2f(4)=-8.2484739\times10^{-9}$, inside the certified enclosure. In the computed $K$, the spread of $K_{st}$ at fixed $s+t-1$ is $2.2\times10^{-19}$ and $|f(d)-f(8-d)|\le1.1\times10^{-19}$, which tests the code against the exact structure.

The full $4\times4$ overlap Gram matrix at $\beta=0$ has two negative eigenvalues, $-6.794442\times10^{-6}$ and $-3.009393\times10^{-6}$ (the largest is $1.510052\times10^{-3}$), about 800 times the $2\times2$ value. The fields $B_1$ and $B_4$ change $U_1(1)$ and $U_1(4)$, so (2) does not hold for this family and the larger witness is a statement at $\beta=0$ only.

{\bfseries\boldmath Positive control.} The Wilson determinant with $r=1$ is link-reflection positive for every real mass [1, 2], and its $\beta=0$ Gram in this geometry must be positive semidefinite. (The theorem is stated in four dimensions; the two-dimensional instance follows by the same computation, which we do not write out.) The same code gives:

\begin{itemize}
\item{} $\det X=\det(D_W-0.7)$, the overlap's own kernel: $f$ decreases toward the antipode ($f(3)-f(4)=+2.6\times10^{-4}$), and the $4\times4$ Gram has eigenvalues $6.1\times10^{-17}$, $1.598\times10^{-2}$, $3.968\times10^{-2}$ and $4.024$, positive semidefinite to rounding.
\item{} $\det(D_W+0.1)$: smallest eigenvalue $-2.1\times10^{-15}$ against a largest of $25.6$, again positive semidefinite to rounding.
\end{itemize}

So the operator builder and the reflection rule give a positive semidefinite Gram where a theorem requires one, and the sign change between $d=3$ and $d=4$ belongs to the overlap.


\section*{5. Discussion}
The failure is in the reflection behaviour of the overlap determinant itself, in a smooth gauge background ($|1-U_p|\le0.1997$ on every plaquette), with a gauge action that is reflection positive. Whether a witness exists inside an admissibility bound smaller than $0.1997$ is not tested here. It does not decide whether a transfer matrix exists for the overlap theory in some other form, whether the violation survives on larger lattices or toward the continuum limit, or whether site reflection fails as well.


\section*{References}
\begin{enumerate}
\item{} K. Osterwalder and E. Seiler, Gauge field theories on a lattice, Ann. Phys. 110 (1978) 440.
\item{} A. Jaffe and B. Janssens, Reflection positive doubles, J. Funct. Anal. 272 (2017) 3506, arXiv:1607.07126.
\item{} H. Neuberger, Exactly massless quarks on the lattice, Phys. Lett. B 417 (1998) 141, arXiv:hep-lat/9707022.
\item{} Y. Kikukawa and K. Usui, Reflection positivity of free overlap fermions, Phys. Rev. D 82 (2010) 114503, arXiv:1005.3751.
\item{} J. E. Mandula, Symmetries of Ginsparg-Wilson chiral fermions, Phys. Rev. D 80 (2009) 085023, arXiv:0901.0572.
\item{} M. Creutz, Positivity and topology in lattice gauge theory, Phys. Rev. D 70 (2004) 091501, arXiv:hep-lat/0409017.
\item{} W. Bietenholz, K. Jansen, K.-I. Nagai, S. Necco, L. Scorzato and S. Shcheredin, Exploring topology conserving gauge actions for lattice QCD, JHEP 03 (2006) 017, arXiv:hep-lat/0511016.
\item{} Y. Kikukawa, On the Hamiltonian formulation of Spin(10) CLGT, talk at the Yukawa Institute for Theoretical Physics, 29 June 2026.
\item{} Z. Lu and S.-H. Shao, On reflection positivity of lattice chiral fermions, arXiv:2611.20084.
\end{enumerate}


\end{document}
